\documentclass[11pt]{article}
\usepackage[margin=1in]{geometry}
\usepackage{booktabs}
\usepackage{array}
\usepackage{tabularx}
\usepackage{xcolor}
\usepackage{hyperref}
\usepackage{enumitem}
\usepackage{titlesec}
\usepackage{parskip}
\hypersetup{colorlinks=true, linkcolor=black, urlcolor=blue!60!black}
\titleformat{\section}{\large\bfseries}{\thesection}{0.6em}{}
\titleformat{\subsection}{\normalsize\bfseries}{\thesubsection}{0.5em}{}

\title{\textbf{Criteria-Augmented Re-scoring of the\\Hide-One-Diagnosis Benchmark}\\[0.5em]
\large 100-patient MIMIC-IV pilot, Haiku predictor}
\author{Thousand Health --- Internal Report}
\date{June 2026}

\begin{document}
\maketitle

\section*{TL;DR}
On a 100-patient MIMIC-IV hide-one-diagnosis pilot using Haiku as the predictor,
strict ICD-label matching scores the system at \textbf{34\% top-3 sensitivity}.
After re-scoring with a 522-entry textbook criteria library (verbatim quotes from
Tier 1--3 sources, re-verified against cited URLs by LLM verifier agents; \emph{not}
physician-reviewed) and an
LLM judge that checks each chart against verbatim guideline thresholds, the
\textbf{criterion-validated sensitivity rises to 78\%} (+44 patients), at a candidate-level
\textbf{lenient PPV of 66\%} (strict 41\%). The criteria layer is doing the real precision
work; Haiku's confidence score adds little additional separation.

\section{Background}
The hide-one-diagnosis benchmark hides a single ICD-coded diagnosis from the chart
and asks the LLM to recover it as one of its top-3 ranked candidates. Strict scoring
requires an exact string/concept match against the hidden ICD label. This scoring
under-counts cases where the LLM names a different but clinically real condition the
patient has --- e.g., the hidden ICD is ``Sialolithiasis'' but the LLM correctly identifies
``Hypokalemia'' that the chart in fact documents.

We hypothesized that many of these ``misses'' are not errors but legitimate co-existing
conditions the strict scorer cannot validate. To test this, we built
\texttt{criteria\_db.json} --- 522 conditions with verbatim quoted textbook criteria from
Tier 1--3 sources (specialty societies, government regulatory, peer-reviewed
foundational papers). Source-verbatim fidelity was checked by LLM verifier agents
that re-fetched each cited URL and confirmed the quoted text appears at the source.
No physician reviewer has audited the library.

\section{Method}
\subsection{Inputs}
\begin{itemize}[leftmargin=*,nosep]
\item \textbf{Predictor}: Claude Haiku, top-3 ranked candidates with a 0--100 confidence
    score per candidate, across 100 MIMIC-IV admissions.
\item \textbf{Criteria DB}: 522 conditions, each with verbatim \texttt{criterion\_quote}
    and source URL. 6 entries dropped during the trust audit for citation errors;
    76 carry caveat notes from the audit.
\item \textbf{Charts}: De-identified MIMIC-IV chart.txt per admission (PHI/DUA-protected;
    not committed).
\end{itemize}

\subsection{Pipeline}
\begin{enumerate}[leftmargin=*,nosep]
\item \textbf{Slug mapping.} LLM mapper assigned each of 218 unique top-3 candidate dx
    strings to a criteria\_db slug or \texttt{null}. 157 mapped (72\%); 61 unmappable
    (vague compound dx, unsupported conditions).
\item \textbf{Criterion-check judging.} For each of 230 (patient, candidate) pairs that
    mapped, a Sonnet judge agent read \texttt{chart.txt}, applied the verbatim criterion
    quote strictly, and returned \texttt{meets} / \texttt{does\_not\_meet} /
    \texttt{insufficient\_evidence}. Ten parallel agents, $\sim$23 jobs each.
\item \textbf{Re-aggregation.} A candidate is a true positive if it either matches the
    strict hidden ICD label \emph{or} the judge said \texttt{meets}. A patient is a
    true positive if any of its top-3 candidates is a TP.
\end{enumerate}

\subsection{Metric definitions}
\begin{description}[leftmargin=*]
\item[Sensitivity (per-patient)] $\text{TP}_{\text{patient}} / 100$, where a patient is a TP
    if $\geq 1$ top-3 candidate above the confidence threshold is validated. Per-patient is
    the only well-defined unit for sensitivity here --- the universe of ``all conditions
    the LLM should have produced'' is unbounded once we leave the strict hidden-ICD scope.
\item[PPV strict (per-candidate)] $\text{TP} / \text{all candidates above threshold}$.
    Counts \texttt{does\_not\_meet}, \texttt{insufficient\_evidence}, and unmappable
    candidates all as false positives. Harsh on the predictor: it pays for our redacted
    labs and incomplete criteria DB.
\item[PPV lenient (per-candidate)] $\text{TP} / (\text{TP} + \text{does\_not\_meet only})$.
    Drops unknowns from the denominator; only counts candidates where the judge actively
    saw chart evidence contradicting the criterion as FP. Fair to the predictor.
\end{description}

\section{Results}

\subsection{Aggregate verdict mix}
Across 230 graded (patient, candidate) pairs:
\begin{center}
\begin{tabular}{lr}
\toprule
Verdict & Count \\
\midrule
meets & 103 \\
does\_not\_meet & 65 \\
insufficient\_evidence & 62 \\
\bottomrule
\end{tabular}
\end{center}

Of the 297 total top-3 candidates: 230 were judgable (slug-mapped), 49 had no slug,
18 had no verdict for other reasons (slug existed but judge had no input row).

\subsection{Patient-level TP rate}
\begin{center}
\begin{tabular}{lrr}
\toprule
Metric & Patients & Rate \\
\midrule
Strict top-3 (hidden ICD recovered) & 34/100 & 34\% \\
Criterion-met top-3 (potential TP only) & 67/100 & 67\% \\
\textbf{Union (strict OR criterion-met)} & \textbf{78/100} & \textbf{78\%} \\
Net new TPs from criteria layer & +44 & \\
\bottomrule
\end{tabular}
\end{center}

Cohort breakdown:
\begin{center}
\begin{tabular}{lrr}
\toprule
Original-scoring cohort & $n$ & Criterion-met in top-3 \\
\midrule
Strict TP (hidden at rank 1--3) & 34 & 23 (sanity check) \\
Matched rank 4--10 (off-podium) & 24 & 16 (66\%) \\
No match in top-10 (rank $-1$) & 42 & 28 (66\%) \\
\bottomrule
\end{tabular}
\end{center}

The 23/34 ``strict TP also passes criteria'' sanity check is reasonable; the other 11
strict TPs had hidden ICDs that did not have a criteria\_db entry, so the criterion
machinery had nothing to validate.

\subsection{Confidence-threshold sweep}
\begin{center}
\small
\begin{tabular}{rrrrrrr}
\toprule
conf $\geq$ & \#cand & TP & PPV strict & PPV lenient & Sens relaxed & Sens strict \\
\midrule
0   & 297 & 121 & 41\% & 66\% & \textbf{78\%} & 34\% \\
60  & 295 & 121 & 41\% & 67\% & 78\% & 34\% \\
65  & 290 & 120 & 41\% & 67\% & 77\% & 34\% \\
70  & 243 & 104 & 43\% & 68\% & 70\% & 31\% \\
75  & 162 &  76 & \textbf{47\%} & \textbf{75\%} & \textbf{56\%} & 24\% \\
80  &  69 &  36 & 52\% & 78\% & 25\% & 15\% \\
85  &  43 &  26 & 60\% & 84\% & 21\% & 13\% \\
90  &  14 &  10 & 71\% & 91\% & 10\% &  7\% \\
\bottomrule
\end{tabular}
\end{center}

\subsection{Confidence calibration}
Mean confidence by ground-truth category:
\begin{center}
\begin{tabular}{lr}
\toprule
Category & Mean confidence \\
\midrule
Strict hidden-match candidates ($n=30$) & 80.7 \\
``meets'' candidates ($n=103$) & 77.7 \\
``does\_not\_meet'' candidates ($n=65$) & 73.6 \\
\bottomrule
\end{tabular}
\end{center}

A 7-point spread between correct and incorrect predictions. Haiku's confidence is only
weakly calibrated --- it is not a reliable filter on its own.

\section{Interpretation}
\subsection{What the +44 actually represents}
The 44 new TPs are cases where Haiku named a different real condition the patient
has --- validated against chart data using verbatim textbook thresholds. Examples from
the ``no-match'' cohort ($n=28$ new TPs):

\begin{itemize}[leftmargin=*,nosep]
\item Hidden: \emph{Sialolithiasis} $\rightarrow$ Haiku ranked Hypokalemia \#1; K was in fact low.
\item Hidden: \emph{Pressure ulcer, stage III} $\rightarrow$ Hypokalemia + Hypophosphatemia, both verified.
\item Hidden: \emph{Melena} $\rightarrow$ Hypomagnesemia + Anemia, both verified.
\item Hidden: \emph{Cervicalgia} $\rightarrow$ Anemia (normocytic) + Hypomagnesemia, both verified.
\end{itemize}

This pattern is loud in the data: ``no-match'' admissions frequently have a relatively
benign hidden ICD (cervicalgia, glaucoma, urticaria), while Haiku correctly surfaces the
clinically significant electrolyte / AKI / anemia / sepsis picture --- which a physician
reading the chart would also prioritize.

\subsection{Where PPV lives}
\begin{itemize}[leftmargin=*,nosep]
\item No threshold: PPV strict 41\%, PPV lenient 66\%, Sens 78\%.
\item conf $\geq 75$: PPV strict 47\%, PPV lenient 75\%, Sens 56\%.
\item conf $\geq 80$: PPV strict 52\%, PPV lenient 78\%, Sens 25\%.
\end{itemize}

The 75 cutoff is the best balanced operating point under Haiku --- modest precision
gain, modest recall cost. Above 80 the sensitivity collapse is non-linear because
Haiku's confidence distribution is compressed (p10 = 68, p90 = 88).

\subsection{Caveats}
\begin{enumerate}[leftmargin=*,nosep]
\item \textbf{Electrolyte / AKI dominance.} $\sim$70\% of the ``meets'' verdicts are
    single-lab-threshold conditions (hypoK, hypoMg, hypoCa, hypoP, AKI by KDIGO).
    These are easy to verify and over-represented in top-3 candidates. ``Real
    test-of-skill'' specific syndromes are a much smaller share.
\item \textbf{62 insufficient\_evidence.} The chart often had redacted lab values
    (\texttt{\_\_\_}). Strict PPV penalizes the predictor for our redactions; the gap
    between strict and lenient PPV is mostly this.
\item \textbf{Slug-mapper was conservative.} 61/218 unique dx strings unmapped. Some
    truly are unmappable (multi-condition free text); some are criteria-DB coverage gaps.
\item \textbf{Predictor is Haiku.} Confidence is poorly calibrated. Re-running on
    Sonnet/Opus would likely shift the curve.
\item \textbf{No physician in the loop, anywhere.} The criteria DB was source-verified
    by LLM agents, not reviewed by clinicians. The 230 chart-vs-criterion verdicts
    were rendered by a single LLM judge (Sonnet), with no physician adjudication.
    Both layers can carry systematic errors a domain expert would catch.
\end{enumerate}

\section{Implications for product framing}
\begin{itemize}[leftmargin=*,nosep]
\item As an \textbf{autonomous flagger} (no physician in loop): 41\% top-3 strict PPV
    is too noisy. A conf-$\geq 90$ slice yields 91\% lenient PPV but 10\% sensitivity ---
    a niche ``high-confidence-only auto-alert'' could be viable.
\item As a \textbf{triage assistant} (physician reviews flagged cases): 78\% relaxed
    sensitivity at 66\% lenient PPV is a defensible operating point. The criteria layer
    transforms the system from a string-match benchmark into a clinically-grounded
    detector.
\item \textbf{The criteria DB is the asset.} Without it, sensitivity is 34\%; with it, 78\%.
    Confidence-thresholding adds little independent lift. The pitch should center on
    verbatim society-guideline criteria as the verification layer, not the LLM's self-reported
    confidence.
\end{itemize}

\section{Next steps}
\begin{enumerate}[leftmargin=*,nosep]
\item Re-run with Sonnet/Opus as predictor; quantify confidence-calibration delta.
\item Expand criteria\_db coverage for unmappable dx strings (61 unique).
\item Spot-check 20--30 ``meets'' verdicts manually for judge-error rate.
\item Replace MDCalc / Wikipedia / weak-domain sources in the 76 caveat entries with
    primary society guideline citations.
\end{enumerate}

\vspace{1em}
\hrule
\vspace{0.5em}
\small
\noindent\textbf{Artifacts.} Re-score detail: \texttt{/tmp/rescore\_result.json} (local,
not committed --- contains PHI-derived data). Judge verdicts: \texttt{/tmp/verdicts\_judge\_0..9.json}.
Criteria DB: \texttt{criteria\_db.json} (522 entries post-audit).

\end{document}
